% Compile with pLaTeX via latexmk; see README.md.
\documentclass[a4paper,12pt,twoside,openany,dvipdfmx]{jsbook}
\usepackage{pethesis}
\usepackage{amsmath}
\usepackage{amssymb}
\usepackage{url}

% Edit all cover information here. Confirm the institutional wording before submission.
\thesistype{令和N年度 修士論文}
\title{情報アクセスシステムの評価\\修士論文の書式確認用サンプル}
\etitle{Evaluating Information Access Systems:\\A Thesis Formatting Demonstration}
\affiliation{早稲田大学大学院 基幹理工学研究科\\情報理工・情報通信専攻}
\supervisor{酒井 哲也 教授}
\researchname{情報アクセス研究}
\studentid{51XXFXXX}
\author{氏名}
\submitdate{202X年X月XX日}

\begin{document}
\maketitle

\begin{coverabstract}
本資料は，日本語の修士論文テンプレートの書式を確認するためのサンプルである．
情報アクセスシステムの評価を題材として，研究背景，関連研究，評価方法，実験結果の説明を例示する．
本文中の数値はすべて書式確認用に作成した仮の値であり，実際の実験結果ではない．
参考文献に含まれる論文の例も架空のものであり，学術的な主張の根拠として利用してはならない．

各章は新しいページから開始するが，奇数ページに限定しない．章の最初のページには下部中央に
ページ番号を表示し，ページ上部の見出しと横線は表示しない．通常の続きのページでは，
偶数ページに章の見出し，奇数ページに節の見出しを表示する．第1章には二つの続きページを設け，
両方のページ上部の表示を比較できるようにした．その他の章にも続きページを設けている．
本文中の明示的な改ページは表示確認のためのものであり，実際の執筆時には取り除いてよい．

このサンプルには目次，数式，表，図，謝辞，参考文献，付録を含める．図目次と表目次は含めない．
表紙と概要にはページ番号を表示せず，目次以降では通し番号を表示する．内容を書き換えた際には，
本文の参照先や図表の配置が正しく更新されていることを，生成したPDFで確認することが望ましい．
\end{coverabstract}

\tableofcontents

% Example files only: remove these inputs and write your thesis directly here.
\chapter{導入}
\label{sec:introduction}

\section{研究の背景}
情報アクセスシステムは，利用者の問いや作業に関連する資料を探すために用いられる．
検索結果を順位付きの一覧として提示する場合もあれば，取得した資料を要約して回答を作る場合もある．
いずれの場合も，何をもって適切な情報と判断するかを明確にしなければ，システムの比較は難しい．
同じ資料であっても，利用者が求める情報の範囲や詳しさによって，役立ち方が異なるためである．

本サンプルでは，検索結果の並べ替えを題材として，修士論文の構成と記述方法を示す．
比較するのは，最初の検索順位をそのまま用いる基準システムと，候補文書に別の得点を与えて
順位を変更するシステムである．ただし，ここに記す手順は説明用の設計例であり，実装や実験を
実施したという意味ではない．実際の論文では，計画と実施済みの内容を区別して記述する必要がある．

\section{研究課題}
検討する課題は二つある．第一に，追加の並べ替えによって上位の検索結果が改善するかを調べる．
第二に，改善が特定の検索要求に偏っていないかを確認する．平均値が上昇したとしても，
すべての検索要求で同じ傾向が得られるとは限らない．一部の要求で大きな改善があり，
別の要求では悪化するという組合せでも，平均値は上昇し得る．

この違いを確認するため，比較では検索要求ごとの得点を保存する．文書集合，検索要求，
適合性判定を両方のシステムで共通にし，変更点を順位付けの処理に限定する．
これにより比較の条件は整理しやすくなるが，評価資料の不足や要求の選び方による不確実性が
なくなるわけではない．結果を解釈する際には，そのような制約も合わせて検討する．

% Demonstration-only break: expose a continuation page.
\newpage
\section{対象範囲と前提}
\label{sec:scope}
本例では，固定された文書集合に対して検索を行い，あらかじめ用意された判定を使って結果を評価する．
実験の途中で文書を追加したり，片方のシステムだけ異なる判定を使用したりしない．
同一条件での比較を行うためには，使用した資料の名称だけでなく，版や更新日も記録することが重要である．
同じ名称の資料であっても，内容が更新されていれば結果が一致しない可能性がある．

評価では，平均的な性能と個々の検索要求の挙動を分けて確認する．平均値は全体を簡潔に示せる一方，
得点の分布を直接表すものではない．個別の結果を保存しておけば，平均値が変化した理由を追跡し，
入力の欠落や識別子の不一致なども確認できる．このような記録は，結果が期待に沿わなかった場合にも必要となる．

\subsection{比較の単位}
比較の単位は一つの検索要求とする．各要求に対して両方のシステムの得点を求め，
その差を記録する．一方の結果が欠けている場合には，単にその要求を取り除くのではなく，
欠落した理由と評価上の取扱いを明記する．除外によって比較対象が変われば，得られた平均値の意味も変わるためである．

\subsection{主張の範囲}
固定された資料を用いた評価から，実際の利用場面での使いやすさを直接結論付けることはできない．
画面の構成，応答時間，利用者の知識など，別の要因も情報の利用に影響する．
本例で扱うのは順位付けの比較手順であり，人を対象とした実験の結果ではない．
結論では，どの条件の下で得られた知見なのかを改めて示す必要がある．

% A second continuation page exposes the other odd/even header position.
\newpage
\section{論文の構成}
第\ref{sec:related_work}章では，関連研究を整理する際の観点を示す．
第\ref{sec:experiments}章では，比較条件，集計方法，表や図の記載例を説明する．
第\ref{sec:approach}章では，評価の問題点と，それに対応する手順の改善案を整理する．
付録\ref{app:details}には，再現に必要な情報と，結果を点検するための記録例を示す．
参照先の番号は自動的に更新されるため，章を追加した場合にも番号を手入力する必要はない．

各章の役割は研究内容に応じて調整してよい．例えば，実装上の工夫が主要な貢献であれば，
システムの構成を説明する独立した章を設ける方法がある．複数の実験が異なる研究課題に対応する場合には，
課題ごとに方法と結果をまとめる構成も考えられる．テンプレートの章立てをそのまま維持することよりも，
問いと証拠の関係を読み手が追えることが重要である．

\section{書式の確認方法}
このページでは，ページ上部に章または節の見出し，下部中央にページ番号が表示される．
偶数ページには章の見出しを左側に，奇数ページには節の見出しを右側に表示する．
一方，章の最初のページでは，上部の見出しと横線を表示せず，ページ番号のみを残す．
この章を三ページ構成とすることで，両方の通常ページの表示を確認できるようにしている．

次の章は，この章の直後のページから始まる．開始ページが偶数であっても空白ページを挿入しない．
サンプル内の明示的な改ページは通常ページの表示を確認するために設けたものである．
実際の執筆時にはこれらを取り除き，本文を自然に流してよい．章の開始時の改ページは，
章を指定する命令によって自動的に行われる．

\chapter{関連研究}
\label{sec:related_work}

\section{文献を整理する観点}
関連研究の章では，既存研究の内容を列挙するだけでなく，自分の研究でどのような判断を行うかに
結び付けて説明する．例えば，検索手法，評価資料，分析方法という観点に分ければ，
採用する基準システムや評価条件を選んだ理由を述べやすくなる．研究課題ごとに整理する場合には，
各課題について何が分かっており，何が未検討であるかを示すとよい．

文献を紹介する際には，対象とする問題，使用した証拠，適用条件を区別する．
同じ手法名が使われていても，対象となる文書や検索要求が異なれば，結果の意味は変わり得る．
そのため，数値だけを抜き出して優劣を判断するのではなく，測定した条件も合わせて記載する必要がある．
これは，後で自分の実験との共通点や相違点を説明するためにも役立つ．

\section{引用の例}
研究室のウェブサイトの項目\cite{RSL}は，参考文献にURLを表示する例として使用している．
日本語論文の項目\cite{exampleArticle}と英語の会議論文の項目\cite{exampleConference}は，
参考文献の書式を確認するための架空の文献である．これらは本文の方法論的な説明を裏付ける資料ではない．
実際の執筆時には，内容を確認した実在の文献に置き換え，引用した文と出典の対応を確かめること．

引用は，その文献がどの説明を支えているかが分かる位置に置く．複数の研究をまとめて述べる場合でも，
すべての文献が同じ条件や結論を扱っているとは限らない．重要な違いがある場合には，
それを文章で説明する．また，書誌情報の正確さと，本文の説明が原文に忠実であることは別の確認事項であり，
一方が正しくても他方が保証されるわけではない．

% Demonstration-only break.
\newpage
\section{比較可能性の確認}
異なる研究の得点を比較する前に，文書集合，要求の選定方法，判定の範囲，評価指標を確認する．
例えば，短い事実確認型の要求を集めた資料と，広い情報収集を目的とする要求を集めた資料では，
上位の結果に求められる内容が異なる可能性がある．得点の差が手法だけによって生じたと考えることはできない．

比較表は設定の違いを簡潔に示すのに適しているが，表だけで解釈を終えてはいけない．
その違いが自分の研究課題にどのように関係するかを説明する必要がある．
従来研究の制約を指摘する場合にも，その研究が本来何を明らかにするために設計されたかを考慮する．
異なる目的に沿った設計を，単純に欠点として扱わないことが重要である．

\section{本研究との関係}
本例では，追加の順位付け処理による変化を，検索要求ごとに対応を付けて調べる．
関連研究の整理は，この手順を選ぶ理由と，検討すべき代替案を明確にする役割を持つ．
基準システムは，単に実装しやすいという理由だけでなく，比較の目的に合っているかという観点から選ぶ．
設定や必要なデータを再現できることも，実験を解釈する上で重要な条件である．

次の章では，これらの観点を具体的な比較条件に落とし込む．
使用する入力をそろえ，変更する部分を限定し，評価結果を集計する方法を先に定める．
関連研究と実験の説明をこのようにつなぐことで，文献調査が実際の研究設計にどのような影響を与えたかを示せる．
なお，以下の条件や数値は，いずれも本テンプレートの表示確認のための例である．

\chapter{評価実験}
\label{sec:experiments}

\section{実験条件の例}
表\ref{tab:settings}は，比較条件の記載方法を示すために作成したものである．
実際にこの条件で実験を実施したわけではない．論文では，使用した設定と，その設定を選択した理由を示す．
結果を見た後で条件を変更した場合には，最初の計画との差異を説明し，追加の分析と区別する必要がある．

\begin{table}[htbp]
  \centering
  \caption{書式確認用の仮の比較条件}
  \label{tab:settings}
  \begin{tabular}{ll}\hline
    項目 & 設定例\\\hline
    候補文書数 & 各検索要求につき100件\\
    評価対象 & 上位10件の検索結果\\
    比較する構成 & 基準システムと並べ替えを行う構成\\
    集計単位 & 検索要求ごとの得点\\
    集計方法 & 各検索要求を等しい重みで平均\\\hline
  \end{tabular}
\end{table}

\section{評価値の集計}
検索要求$i$に対する評価値を$s_i$，比較対象の要求数を$n$とすると，平均値は次式で表される．
\begin{equation}
  \label{eq:mean}
  \bar{s}=\frac{1}{n}\sum_{i=1}^{n}s_i.
\end{equation}
式\eqref{eq:mean}は，各要求に同じ重みを与える集計である．
この式だけでは，適合性や順位をどのように評価値へ変換するかは定まらない．
実際には，$s_i$を求める指標の定義や，未判定文書の扱いも明示する必要がある．

実験の記録には，要求識別子と両システムの得点を残し，設定ファイルや実行ログと対応付ける．
この記録方法の改善案を次章で説明する．

% Demonstration-only break; flush the tables before showing a continuation page.
\clearpage
\section{結果と図の配置例}
表\ref{tab:results}の値は，桁の配置と参照の動作を確認するための仮の値である．
二つの群の要求数が等しいという設定の下で，全体の平均を示している．
この表から，実在のシステムの有効性や統計的な確かさを主張することはできない．

\begin{table}[htbp]
  \centering
  \caption{実験結果ではない仮の評価値}
  \label{tab:results}
  \begin{tabular}{lrrr}\hline
    検索要求の群 & 基準 & 比較対象 & 差\\\hline
    短い要求 & 0.420 & 0.460 & +0.040\\
    詳細な要求 & 0.380 & 0.390 & +0.010\\
    全体 & 0.400 & 0.425 & +0.025\\\hline
  \end{tabular}
\end{table}

図\ref{fig:logo}は，図の中央配置，図題，本文からの参照を確認するための例である．
元のテンプレートに含まれていた校徽を使用しており，研究結果を示すものではない．
実際の論文では，図から読み取るべき内容と，その内容が研究課題にどう関係するかを本文で説明する．

\insertfigure{waseda_logo.pdf}{書式確認用に掲載した早稲田大学の校徽}{fig:logo}{width=3cm}

\section{結果を説明する際の注意}
平均値の上昇は，すべての要求で改善したことを意味しない．
個別の得点を調べ，改善と悪化がどのように分布しているかを確認する必要がある．
大きく変化した要求を取り上げる場合には，選択した基準を記述し，印象的な例だけで全体の傾向を
説明しようとしないことが重要である．具体例は挙動を理解する助けになるが，発生頻度の推定とは異なる．

考察では，観察した事実と，その原因についての仮説を分ける．
例えば，ある要求で無関係な文書が上位に現れた場合，語の曖昧さが関係している可能性はある．
しかし，一つの結果だけで一般的な原因を確定することはできない．
追加の事例や条件を制御した実験によって，その解釈を支える証拠を集める必要がある．


\chapter{問題点と解決提案手法}
\label{sec:approach}

\section{評価手順に残る問題}
評価資料には，判定が十分に付与されていない文書や，特定の分野に偏った検索要求が含まれる場合がある．
そのような資料を用いた比較では，得点が何を表しているかを慎重に解釈する必要がある．
手法の有効性と，資料の構成によって生じる差を完全に切り分けることは容易ではない．
まず，確認できた制約を記録し，結論の範囲を明確にすることが重要である．

実装上の問題も比較の信頼性に影響する．入力の欠落，要求識別子の重複，文書の正規化処理の違いなどがあると，
意図した条件で比較できなくなる．これらの問題は，結果の平均値だけを見ても発見しにくい．
評価前に入力と出力の対応を確認する処理を設けることで，不完全な結果をそのまま集計することを防げる．

\section{提案する確認手順}
本例で提案するのは，検索要求ごとの対応関係を保存し，集計前に欠落と重複を点検する手順である．
両方の結果がそろっていることを確認した後に差を求め，群ごとの集計へ進む．
欠落がある場合には，その扱いを評価規則として明示する．自動的に得点をゼロに置き換える処理は，
その規則が意図したものか確認した上で採用する必要がある．

この手順では，比較条件と評価結果を別々のファイルとして保存し，実行識別子で結び付ける．
条件は何を行う予定だったかを示し，実行ログは実際に何が起きたかを示す．
予想外の挙動があった場合にも，その情報を捨てずに残すことが，再現と原因の調査に役立つ．
ここで述べた手順は説明用の提案であり，効果を実証した結果ではない．

% Demonstration-only break.
\newpage
\section{検証方法}
提案した確認手順自体も点検する必要がある．少数の検索要求を用いた例で，
入力識別子の対応，得点の計算，平均値への集計を手作業で追えるようにする．
その上で，欠落や重複を含む入力に対して意図した動作になるかを確認する．
正常な入力で動くだけでは，誤った入力を検出できることまでは保証されない．

処理の再現性を確認する際には，乱数の種だけでなく，ソフトウェアと資料の版も記録する．
並列処理や外部サービスに依存する部分がある場合には，完全に同じ結果が得られない要因を説明する．
再実行で変化する範囲が分かれば，その変化と手法による差を区別して検討しやすくなる．

\section{今後の課題}
次の段階として，異なる特徴を持つ検索要求の集合で，同じ比較を実施することが考えられる．
追加する資料は，単に規模を増やすためではなく，現在の解釈を制約している不確実性を調べるために選ぶ．
元の実験と共通にする条件と，変更が必要な条件を区別すれば，得られた違いを説明しやすくなる．

計算時間やメモリ使用量も検討対象となり得る．順位の改善が得られたとしても，
応答時間の増加が利用場面に適しているとは限らない．測定に含める処理の範囲を定め，
実行環境とともに結果を報告する必要がある．どの課題から取り組むかは，
研究の主張を支える上で最も重要な未解決点が何かに基づいて判断する．


\begin{acknowledgement}
ここには実際の研究に対する指導，助言，技術的支援などについて謝辞を記載する．
この段落は書式確認のための例文であり，実際の支援内容や研究費の情報を記述したものではない．
論文として提出する前に，自身の研究で受けた支援に即した文章に置き換えること．

謝辞は章番号を付けないが，目次には掲載する．このページでも，通常の章の最初のページと同様に，
下部中央のページ番号だけを表示する．指導者や協力者に言及する際には，その人が論文の結論を
承認したと誤解されるような表現を避け，どのような協力を受けたかを具体的に記述するとよい．
\end{acknowledgement}

\bibliographystyle{unsrt}
\bibliography{reference}

% Remove this appendix and its input if your thesis does not need an appendix.
\appendix
\chapter{実験結果詳細}
\label{app:details}

\section{再現のための記録}
付録では，本文の流れを妨げずに詳細な情報を残すことができる．
本例では，実験を再現するために記録しておく項目を整理する．
資料の名称だけでなく，実際に使用した版や設定を記載し，本文中の表と結果ファイルの対応が
分かるようにすることが望ましい．ファイルを保存するだけでなく，どの処理で作られたかも説明する．

確認する項目の例を以下に示す．
\begin{enumerate}
  \item 文書集合，検索要求，適合性判定の版を記録する．
  \item 前処理と順位付けに使用した設定を保存する．
  \item 検索要求ごとの出力と評価値を対応付ける．
  \item 除外した要求とその理由を明記する．
  \item 各表を再生成するための集計手順を記載する．
\end{enumerate}

付録の内容は，本文から参照されることで役割が明確になる．
例えば，第\ref{sec:experiments}章の集計方法を再現するために必要な設定を示す，という関係を説明するとよい．
ページ番号を手入力する代わりに参照命令を使用すれば，本文が増減しても参照先を更新できる．
このサンプルでは，付録にアルファベットの番号を付け，節番号にも同じ文字を使用する．

% Demonstration-only break.
\newpage
\section{文書全体の点検}
このページは，付録の通常ページの書式を確認するための例である．
付録の最初のページにはページ番号だけを表示し，続きのページには見出しと横線も表示する．
目次に示される付録のページ番号は，付録の最初のページを指していなければならない．
本文の各章と同様に，付録の前にも奇数ページへ合わせるための空白ページは挿入しない．

図や表を追加すると，本文のページ分割が変化することがある．
これは通常の組版の動作であるが，図題と図が離れていないか，参照先の番号が更新されたか，
大きな表が余白を越えていないかを確認する必要がある．長い見出しについても，
本文とページ上部の見出しの両方で，文字が重なっていないかを点検する．

正式な論文を執筆する際には，例文と架空の数値，架空の文献を削除する．
表示確認のために設けた改ページも取り除き，実際の内容を主ファイルに記述してよい．
表紙の所属，学籍番号，氏名，提出日は実際の情報に置き換える．
文書の体裁が整っていても，内容の正確さが自動的に保証されるわけではないため，両方を確認することが重要である．

\end{document}
